% ========================
% 3.5 Implementierungsdetails
% ========================
\section{Implementierungsdetails}


% ========================
% 3.5.1 Datenmodell
% ========================
\subsection{Datenmodell}
\setauthor{\firstauthor}

Das Datenmodell bildet die grundlegende Struktur der gesamten Anwendung und wurde als gemeinsames Schema 
von beiden Entwicklungsteams entworfen. Es umfasst alle zentralen Entitäten zur Verwaltung von Nutzer:innen,
Fragen, Themenbereichen, Übungsfortschritt und Prüfungssimulation.

\begin{figure}[H]
  \centering
  \fbox{
    \includegraphics[width=\textwidth]{datamodel.png}
  }
  \caption{Datenbankmodell der Anwendung}
  \label{fig:datamodel}
\end{figure}

\autoref{fig:datamodel} zeigt das vollständige relationale Datenbankmodell der Anwendung. Nicht alle 
dargestellten Tabellen werden in jedem beschriebenen Use Case dieser Arbeit verwendet, da einzelne Komponenten 
von unterschiedlichen Teams umgesetzt wurden. Für die folgenden Abschnitte werden daher nur jene Entitäten 
näher betrachtet, die für die Umsetzung der Übungs- und Prüfungsfunktionen relevant sind.


% ========================
% 3.5.1.1 Übungsmodus
% ========================
\subsubsection{Zentrale Entitäten für den Übungsmodus}
Für den Übungsmodus sind insbesondere die Entitäten \textit{Question}, \textit{Answer}, \textit{QuestionPool}, 
\textit{QuestionPoolEntry} und \textit{User} relevant. Sie bilden die Grundlage, um Fragen darzustellen, Antworten 
zu erfassen und den Lernfortschritt der Nutzer:innen nachzuvollziehen.

\vspace{0.2cm}
\textbf{Question} \\
Die Entität \textit{Question} repräsentiert eine einzelne Übungsfrage. Sie enthält den Fragetext und den Fragetyp
(\textit{MULTIPLE CHOICE}, \textit{TEXT}). Jede Frage kann mehrere Antworten besitzen, die in der Entität \textit{Answer} 
hinterlegt sind.

\vspace{0.2cm}
\textbf{Answer} \\
\textit{Answer}-Objekte beinhalten den Antworttext und einen Eintrag, ob die Antwort korrekt ist (\textit{isCorrect}). 
Unabhängig vom Fragetypen werden alle richtigen und falschen Antworten auf eine Frage gespeichert. Bei Multiple-Choice-Fragen 
müssen alle richtigen und keine falschen Antworten angegeben werden, um als korrekt gewertet zu werden. Für Freitext-Fragen 
werden nur richtige Antworten angelegt. Die eingegebene Antwort muss hierbei einer der Einträge gleichen, um korrekt gewertet
zu werden.

\vspace{0.2cm}
\textbf{QuestionPool und QuestionPoolEntry} \\
Der individuelle Fragenpool der Nutzer:innen wird von einem \textit{Questionpool}-Objekt repräsentiert. Die Entität 
\textit{QuestionPoolEntry} verbindet eine Frage mit einem Pool und speichert den individuellen Fortschritt des Nutzers/der Nutzerin. 
Dazu zählen unter anderem:

\begin{itemize}
  \item \textit{correctCount}: Anzahl der aufeinanderfolgend korrekt geantworten Male für diese Frage
  \item \textit{lastAnsweredCorrectly}: ob die Frage zuletzt richtig beantwortet wurde
  \item \textit{answeredAt}: Zeitpunkt, zu dem die Frage zuletzt beantwortet wurde
\end{itemize}

Diese Tabelle ermöglicht die Umsetzung der Wiederholungslogik im Übungsmodus, indem Fragen je nach Bearbeitungsstatus 
und Lernerfolg wiederholt oder vorübergehend gesperrt werden.

\vspace{0.2cm}
\textbf{User} \\
Die Entität \textit{User} enthält die grundlegenden Informationen über die Lernenden. Sie ist mit den \textit{QuestionPool}-Objekten 
verbunden, um den individuellen Lernfortschritt zu speichern, sowie mit \textit{Question} und \textit{Solution}, um erstellte Fragen und 
Lösungswege zu verwalten.

\vspace{0.2cm}
\textbf{Funktionale Zusammenhänge} \\
In \autoref{fig:datamodel} sind die Beziehungen zwischen diesen Entitäten dargestellt. Zusammen ermöglichen sie:

\begin{itemize}
  \item Die flexible Auswahl von Fragen aus einem Pool für eine Übungssitzung
  \item Das Tracking des individuellen Lernfortschritts pro Frage
  \item Die Umsetzung von Wiederholungs- und Sperrlogiken im Übungsmodus
\end{itemize}

Diese Struktur ermöglicht, dass im Übungsmodus flexibel auf die Bedürfnisse der Nutzer:innen eingangen werden kann, während
gleichzeitig eine saubere Trennung zwischen Fragen, Fragenpools der Nutzer:innen und  individuellen Fortschritten besteht.


% ========================
% 3.5.1.2 Lösungswege
% ========================
\subsubsection{Zentrale Entitäten für Lösungswege}
Für die Darstellung von Lösungswegen im Übungsmodus sind vor allem die Entitäten \textit{Solution}, \textit{SolutionStep} und 
\textit{SolutionVote} relevant. Durch diese können Nutzer:innen nicht nur die richtige Lösung sehen, sondern auch die einzelnen Schritte 
nachvollziehen und Feedback von anderen Nutzer:innen erhalten. Zudem können individuelle Lösungswege hochgeladen werden, um 
sich und anderen Nutzer:innen in der Zukunft zu unterstützen.

\vspace{0.2cm}
\textbf{Solution} \\
Die Entität \textit{Solution} repräsentiert einen Lösungsweg zu einer Frage. Jede Lösung gehört zu genau einer Frage (\textit{Question}) 
und einem erstellenden Nutzer (\textit{User}) und stellt sich aus mehreren \textit{SolutionStep}-Objekten zusammen

\vspace{0.2cm}
\textbf{SolutionStep} \\
\textit{SolutionStep}-Objekte enthalten die einzelnen Schritte einer Lösung. Jeder Schritt besitzt einen Titel zur Zusammenfassung und 
einen Text zur Erklärung und Veranschaulichung des nächsten Schritts. Durch die Schritt-für-Schritt-Darstellung kann der/die Nutzer:in 
den Lösungsweg nachvollziehen und gezielt lernen.

\vspace{0.2cm}
\textbf{SolutionVote} \\
Diese Entität ermöglicht es anderen Nutzer:innen, eine Lösung positiv oder negativ zu bewerten (\textit{isUpVote}). So werden Lösungswege 
hochwertige Lösungswege mit mehr positiven Bewertungen hervorgehoben und die Community kann gegenseitig Feedback geben. 

\vspace{0.2cm}
\textbf{Funktionale Zusammenhänge} \\

\begin{itemize}
  \item Jede Frage kann mehrere von Nutzer:innen erstellte Lösungen besitzen.
  \item Jede Lösung besteht aus mehreren Schritten.
  \item Nutzer:innen können Lösungswege bewerten, wodurch die Qualität sichtbar wird.
  \item Die Anzeige der Lösungen erfolgt im Übungsmodus direkt nach der Beantwortung, wodurch der Lernprozess unterstützt wird.
\end{itemize}



% ========================
% 3.5.1.3 Prüfungsmodus
% ========================
\subsubsection{Datenmodell für Prüfungssimulationen}
Für die Umsetzung des Prüfungsmodus sind relevante Entitäten vor allem \textit{Exam} und \textit{ExamQuestion}. Durch diese sind die
Verwaltung von Prüfungen, die Speicherung der Antworten und die nachträgliche Auswertung möglich.

\vspace{0.2cm}
\textbf{Exam} \\
Die Entität \textit{Exam} stellt eine einzelne Prüfung dar. Sie enthält Informationen über das Zeitlimit, den Beginnzeitpunkt, den 
Endzeitpunkt sowie dem Gesamtergebnis. Über die Beziehung zu \textit{User} ist klar, welcher Nutzer/welche Nutzerin diese Prüfung
abgelegt hat. Jede Prüfung kann mehrere Fragen enthalten, die über die Entität \textit{ExamQuestion} abgebildet werden.

\vspace{0.2cm}
\textbf{ExamQuestion} \\
\textit{ExamQuestion} speichert die gegebene Antworten auf die einzelne Frage innerhalb einer Prüfung und ist mit dem dementsprechendem
Fragenpoolobjekt verknüpft. Sie speichert nicht nur die ausgewählte Antwort, sondern auch ob diese Frage korrekt beantwortet wurde. So
können alle Fragen einer Prüfung individuell ausgewertet werden und im Reviewmodus angezeigt werden.

\vspace{0.2cm}
\textbf{Funktionale Zusammenhänge} \\
Durch die Beziehung zwischen \textit{Exam} und \textit{ExamQuestion} wird folgendes ermöglicht:

\begin{itemize}
  \item Speicherung von Prüfungen für einzelne Nutzer:innen
  \item Zuordnung der ausgewählten Antworten zu den jeweiligen Fragen
  \item Berechnung des Gesamtergebnisses und der prozentual korrekt beantworteten Fragen
  \item Nachträgliche Ansicht der Prüfung im Reviewmodus, inklusive der Anzeige von korrekt und falsch beantworteten Fragen
\end{itemize}

Diese Struktur macht eine realistische Simulation einer Prüfungssituation möglich. Die Nutzer:innen bearbeiten Fragen unter Zeitdruck, 
ihre Antworten werden erfasst und persisitiert und eine nachträgliche Auswertung ist jederzeit möglich.



% ========================
% 3.5.2 Überblick über Komponenten
% ========================
\subsection{Überblick über Komponenten}
\setauthor{\firstauthor}


% ========================
% 3.5.2.1 Backend
% ========================
\subsubsection{Backend / API}
Das auf Quarkus basierende Backend stellt über REST-Endpunkte die zentrale Schnittstelle für das Frontend bereit. Dabei werden alle Daten
über JSON-Format übertragen, wobei klassische HTTP-Methoden wie \textit{GET} und \textit{POST} verwendet werden.

Die Kernaufgaben des Backends umfassen:

\begin{itemize}
  \item \textbf{Fragenmanagement:} \\
  Laden von Fragen, Abrufen von Antwortoptionen und Speichern der Nutzerantworten. Die Logik zur Bewertung der Antworten wird durch Services 
  wie \textit{AnswerService} umgesetzt. 

  \item \textbf{Prüfungslogik:} \\
  Erstellung und Verwaltung von Prüfungen über den \textit{ExamService}. Dieser Service wählt Fragen aus den individuellen Fragenpools aus, 
  setzt Zeitlimits und berechnet den Endpunktestand. Auch Funktionen wie das Kopieren bereits abgeschlossener Prüfungen werden hier ermöglicht.

  \item \textbf{Fortschritts- und Statistikfunktionen:} \\
  Der \textit{StatsService} liefert alle Daten für die Fortschritts- und Statistikübersichten im Frontend. Dazu gehören globale Statistiken, 
  Fortschritte in einzelnen Themenbereichen sowie die Umsetzung der Sperrlogik für Übungsfragen, um Wiederholungen strukturiert zu steuern.

  \item \textbf{Kommunikation mit dem Frontend:} \\
  Das Backend stellt alle relevanten Daten für den \textit{Question Runner}, die Fortschrittsübersicht und die Ergebnisanzeige der Prüfungen 
  bereit. Eingehende Antworten werden validiert, ausgewertet und in der Datenbank persistiert.

  \item \textbf{Datenzugriff:} \\
  Services greifen auf die Daten über dedizierte Repository-Klassen (wie zum Beispiel \textit{QuestionRepository} oder \textit{ExamRepository}) 
  zu. Das führt zu einer klaren Trennung zwischen Geschäftslogik und Datenzugriff.
\end{itemize}

Das Backend ist unter anderem für besondere Funktionen, wie die Wiederholungslogik für Übungsfragen, die Sicherstellung der Prüfungskonsistenz 
sowie die Kopie von abgeschlossenen Prüfungen, verantwortlich. Dadurch bleibt die Trennung zwischen Übungs- und Prüfungsmodus sauber erhalten
und im Frontend kann die Darstellung vollständig im Fokus stehen.

\begin{figure}[H]
  \centering
  \fbox{
    \includegraphics[width=0.7\textwidth]{component-diagramme-backend.png}
  }
  \caption{Komponentendiagramm des Backends}
\end{figure}


% ========================
% 3.5.2.2 Frontend
% ========================
\subsubsection{Frontend-Komponenten}
Das Frontend der Anwendung ist sowohl für die Darstellung der Inhalte als auch für die Interaktion der Nutzer:innen mit dem Übungs- und 
Prüfungsmodus verantwortlich. Dabei ist eine klare Struktur und Wiederverwendung zentraler Komponenten Grundlage für die Gewährleistung 
von konsistenten Nutzererlebnissen über alle Modi hinweg.

Eine zentrale Rolle nimmt dafür der \textbf{Question Runner} ein. Diese Komponente wird sowohl im Übungs-, Prüfungs- als auch im Reviewmodus
verwendet und ist für die Darstellung einzelner Fragen, der zugehörigen Antwortmöglichkeiten sowie der Navigation zwischen Fragen 
verantwortlich. Je nachdem, welcher Modus aktiv ist, verändert sich dabei das Verhalten der Komponente, in Bezug auf Feedback, 
Darstellung des Zeitlimits oder Bearbeitbarkeit von Antworten, während die grundlegende Oberfläche unverändert bleibt.

\begin{figure}[H]
  \centering
  \fbox{
    \includegraphics[width=0.7\textwidth]{component-diagramme-frontend.png}
  }
  \caption{Komponentendiagramm des Frontends}
\end{figure}

In die jeweiligen Modi kann über eigene Views eingestiegen werden, wie etwa die Fortschrittsübersicht im Übungsmodus oder die Prüfungskonfiguration
im Prüfungsmodus. Diese Komponenten sind vor allem für die Auswahl von Themenbereichen und den Start eines Durchlaufs zuständig und übergeben
lediglich die notwendigen Parameter an den Question Runner.

Diese Aufteilung gewährleistet eine klare Trennung zwischen Darstellungslogik und fachlicher Logik. Zur gleichen Zeit ermöglicht dieser 
modulare Aufbau eine einfache Erweiterbarkeit für weitere Modi, ohne bestehende Komponenten grundlegend anpassen zu müssen.


% ========================
% 3.5.3 Custom CSS-Klassen und eingebundene Bibliotheken
% ========================
\subsection{Custom CSS-Klassen und eingebundene Bibliotheken}
\setauthor{\secondauthor}

Die Implementierung der Lernplattform erfordert die Integration externer Bibliotheken zur Datenvisualisierung sowie wiederverwendbare CSS-Komponenten für ein konsistentes Erscheinungsbild.

\subsubsection{Visualisierung von Statistiken mit Chart.js}

Um den Lernfortschritt der Nutzer:innen übersichtlich darzustellen, wird Chart.js verwendet. Die Open-Source-Bibliothek integriert sich nahtlos in Angular und bietet reaktive Donutdiagramme mit dynamischen Datenaktualisierungen. Chart.js wurde gewählt wegen seiner Performance, umfangreichen Dokumentation und der \texttt{update()}-Methode für flüssige Animationen ohne Seitenneuladen.

\textbf{Donutdiagramm-Implementierung}

Donutdiagramme zeigen prozentuale Verteilungen des Lernkarten-Status (offen, 1x richtig, 2x richtig, 3x richtig, falsch). Der zentrale Freiraum enthält Schlüsselkennzahlen (siehe \autoref{fig:chartjs}).

\begin{figure}[H]
    \centering
    \fbox{\includegraphics[width=0.7\textwidth]{chartjs.png}}
    \caption{Donutdiagramm: Globale Lernstatistik mit Farbkodierung}
    \label{fig:chartjs}
\end{figure}

\textbf{Datenfluss: Von der Datenbank zum Canvas-Rendering}

\begin{enumerate}
\item \textbf{Angular Lifecycle: ngOnInit()}\\
Wenn die Komponente geladen wird, startet sie das Laden der Statistikdaten. Der \texttt{StatsService} wird dafür verwendet, damit die Komponente selbst nicht direkt mit dem Backend arbeiten muss.
\begin{verbatim}
ngOnInit(): void {
  this.statsService.getStatsOverview(this.userId).subscribe(data => {
\end{verbatim}

\item \textbf{REST API Request: HttpClient}\\
Angular sendet eine GET-Anfrage an das Backend und übergibt dabei die Benutzer-ID in der URL.

\item \textbf{JAX-RS Controller}\\
Im Backend wird diese Anfrage von der passenden Methode angenommen und an den Statistik-Service weitergeleitet.
\begin{verbatim}
@GET @Path("/{userId}/overview")
public StatsOverviewDto getStatsOverview(Long userId) {
  return statsService.calculateStatsOverview(userId);
}
\end{verbatim}

\item \textbf{JPA Repository: SQL Aggregation}\\
Im Hintergrund wird eine SQL-Abfrage ausgeführt, die die Antwort-Status des Benutzers zusammenzählt. Dadurch muss nicht jede Frage einzeln geladen werden.
\begin{verbatim}
SELECT COUNT(CASE WHEN status='UNANSWERED' THEN 1 END) as unanswered,
       COUNT(CASE WHEN status='CORRECT_ONCE' THEN 1 END) as correctOnce
FROM question_pool_entries WHERE user_id = 123;
\end{verbatim}

\item \textbf{JSON Serialisierung: Jackson}\\
Das Backend gibt die berechneten Werte als JSON zurück, damit sie im Frontend weiterverwendet werden können.
\begin{verbatim}
{
  "unanswered": 15, "correctOnce": 25,
  "legend": [{"label":"Offen","value":15,"color":"#E0E0E0"}]
}
\end{verbatim}

\item \textbf{Data-Binding: RxJS und Angular Change Detection}\\
Die Antwort wird im Frontend entgegengenommen und in das Format umgewandelt, das Chart.js benötigt. Hier werden vor allem Labels, Werte und Farben aus dem Backend übernommen.
\begin{verbatim}
this.doughnutChartData = {
  labels: data.legend.map(e => e.label),
  datasets: [{data: data.legend.map(e => e.value),
              backgroundColor: data.legend.map(e => e.color)}]
};
\end{verbatim}
Wenn sich diese Daten ändern, erkennt Angular das automatisch und bereitet das Diagramm für ein neues Rendering vor.

\item \textbf{Canvas Rendering: Chart.js und ng2-charts}\\
Chart.js zeichnet das Donutdiagramm auf ein HTML-Canvas. Die Angular-Direktive \texttt{ng2-charts} verbindet dabei die Angular-Daten mit Chart.js.
\begin{verbatim}
<canvas baseChart [data]="doughnutChartData" [type]="doughnutChartType">
\end{verbatim}
Sobald sich \texttt{doughnutChartData} ändert, wird das Diagramm automatisch aktualisiert und neu dargestellt.
\end{enumerate}

\subsubsection{Wiederverwendbare CSS-Klassen}

Die Anwendung nutzt ein System wiederverwendbarer CSS-Klassen, um ein konsistentes Design und Layout über alle Komponenten hinweg zu gewährleisten. Diese Klassen sind in einer zentralen, geteilten CSS-Datei definiert und können komponentenübergreifend verwendet werden. Dieser Ansatz folgt dem DRY-Prinzip (Don't Repeat Yourself) und bringt mehrere Vorteile mit sich:

\begin{itemize}
    \item \textbf{Konsistenz}: Ein einheitliches Erscheinungsbild der Anwendung wird durch die Verwendung derselben Stilklassen in verschiedenen Komponenten sichergestellt.
    \item \textbf{Wartbarkeit}: Änderungen am Design müssen nur an einer zentralen Stelle vorgenommen werden und wirken sich automatisch auf alle verwendenden Komponenten aus.
    \item \textbf{Teamarbeit}: Die gemeinsame CSS-Datei vereinfacht die Zusammenarbeit im Team, da alle Entwickler:innen auf denselben Stil-Definitionen aufbauen.
    \item \textbf{Reduzierung von Redundanz}: Durch die Wiederverwendung werden doppelte Code-Definitionen vermieden, was die Dateigröße reduziert und die Ladezeit der Anwendung optimiert.
\end{itemize}

\textbf{Beispiele häufig verwendeter CSS-Klassen}

Die folgenden Klassen bilden das Grundgerüst des visuellen Designs der Lernplattform:

\begin{itemize}
    \item \textbf{.widget-wrapper}: 
    Diese Klasse definiert den Grundcontainer für die Widget-Darstellung der Anwendung. Widgets sind modulare UI-Komponenten, die Informationen kompakt und übersichtlich präsentieren, wie beispielsweise Statistikkarten, Fortschrittsanzeigen oder Lernziel-Übersichten. 
    
    Die \texttt{.widget-wrapper}-Klasse sorgt für ein einheitliches Layout aller Widgets durch standardisierte Abstände, Padding-Werte, Border-Radius und Schatten-Effekte. Dies gewährleistet, dass verschiedene Widget-Typen visuell harmonieren und als zusammengehörige Elemente wahrgenommen werden. Zusätzlich implementiert die Klasse responsive Breakpoints, sodass Widgets auf verschiedenen Bildschirmgrößen optimal dargestellt werden.
    
    \item \textbf{.header}:
    Die Header-Klasse fügt, wo benötigt, einen visuell hervorgehobenen Kopfbereich hinzu, der mit der primären Farbe des Anwendungsdesigns gestaltet ist. Diese Klasse wird typischerweise für Überschriften von Widgets, Sektionen oder modalen Dialogen verwendet.
    
    Die primäre Designfarbe schafft dabei eine klare visuelle Hierarchie, sodass Nutzer:innen wichtige Bereiche der Anwendung schnell erkennen können. Neben der Hintergrundfarbe legt die Klasse auch Textfarbe, Schriftgröße, Schriftgewicht und vertikale Abstände fest, damit alle Header-Elemente einheitlich dargestellt werden.
    
    \item \textbf{button.submit-btn}:
    Diese Klasse stellt einen standardisierten Button-Stil für Interaktionen bereit und kommt vor allem bei Übungsaufgaben, Formular-Submissions und anderen Bestätigungsaktionen zum Einsatz. 
    
    Der \texttt{.submit-btn} zeigt visuelles Feedback für unterschiedliche Zustände: normale Darstellung, Hover-Effekt, aktiver Zustand beim Klicken und deaktivierter Zustand. Nutzer:innen können dadurch schnell erkennen, welche Schaltflächen primäre Aktionen auslösen. Übergänge und Animationen sorgen dafür, dass Interaktionen flüssiger wirken.
    
    Die Button-Klasse berücksichtigt auch Accessibility-Aspekte wie ausreichende Größe für Touch-Bedienung, gute Farbkontraste und sichtbare Fokus-Indikatoren bei Tastatur-Navigation.
\end{itemize}

\textbf{CSS-Architektur und Best Practices}

Die CSS-Klassen sind modular strukturiert und unterscheiden zwischen globalen Basisklassen, komponentenspezifischen Stilen und Utility-Klassen. So lassen sich einzelne Komponenten flexibel anpassen, ohne dass globale Stile beeinflusst werden. Außerdem treten dadurch weniger Spezifität-Konflikte auf, was die Wartbarkeit des Codes verbessert.

